% =========================================================
\section{Semaine 3 : Résolution de systèmes linéaires}
% =========================================================

\subsection{Introduction : Résolution de systèmes linéaires}

\textbf{But :} Résoudre des systèmes de taille $n \times n$, de la forme :
\[ A\vec{x} = \vec{b} \]
où $A = (a_{ij})_{1 \le i,j \le n}$ est une matrice carrée, $\vec{x}$ est le vecteur des inconnues et $\vec{b}$ le vecteur second membre.

\textbf{Exemples d'application :}
\begin{itemize}
    \item Méthode de Newton pour les systèmes non linéaires.
    \item Résolution d'Équations aux Dérivées Partielles (EDP) discrétisées (cf. suite du cours).
\end{itemize}

On sait que si $A$ est inversible, le système admet une unique solution. \\
\textbf{Question :} Comment la trouver efficacement ?

\subsection{Systèmes triangulaires}

\subsubsection{Système triangulaire inférieur}
Un système est triangulaire inférieur si $a_{ij} = 0$ pour tout $j > i$.

\begin{algorithme}[Substitution progressive]
\textbf{Entrées :} Matrice $A \in \mathbb{R}^{n \times n}$ triangulaire inférieure, vecteur $\vec{b} \in \mathbb{R}^n$. \\
\textbf{Pour} $i = 1, \dots, n$ \textbf{faire :}
\[ x_i = \frac{b_i - \sum_{j=1}^{i-1} a_{ij}x_j}{a_{ii}} \]
\textbf{Fin Pour} \\
\textbf{Retourne :} $\vec{x}$, solution de $A\vec{x} = \vec{b}$.
\end{algorithme}

\subsubsection{Système triangulaire supérieur}
Un système est triangulaire supérieur si $a_{ij} = 0$ pour tout $j < i$.

\begin{algorithme}[Substitution rétrograde]
\textbf{Entrées :} Matrice $A \in \mathbb{R}^{n \times n}$ triangulaire supérieure, vecteur $\vec{b} \in \mathbb{R}^n$. \\
\textbf{Pour} $i = n$ à $1$ \textbf{faire :}
\[ x_i = \frac{b_i - \sum_{j=i+1}^{n} a_{ij}x_j}{a_{ii}} \]
\textbf{Fin Pour} \\
\textbf{Retourne :} $\vec{x}$, solution de $A\vec{x} = \vec{b}$.
\end{algorithme}

\subsection{Factorisation LU}

\subsubsection{Principe et élimination de Gauss}
L'objectif est de décomposer la matrice $A$ sous la forme $A = LU$, où $L$ est triangulaire inférieure et $U$ est triangulaire supérieure.

\textbf{Exemple : Résoudre $A\vec{x} = \vec{b}$}
\[ 
A = \begin{pmatrix} 2 & 1 & 0 \\ -4 & 3 & -1 \\ 4 & -3 & 4 \end{pmatrix}, \quad 
\vec{b} = \begin{pmatrix} 4 \\ 0 \\ -2 \end{pmatrix} 
\]
Application de l'élimination de Gauss :
\begin{enumerate}
    \item $L_2 \leftarrow L_2 - \left(\frac{-4}{2}\right)L_1$
    \item $L_3 \leftarrow L_3 - \left(\frac{4}{2}\right)L_1$
\end{enumerate}
On obtient une étape intermédiaire. Puis on annule le coefficient de la 3ème ligne, 2ème colonne :
\begin{enumerate}
    \item $L_3 \leftarrow L_3 - \left(\frac{-5}{5}\right)L_2$
\end{enumerate}
À la fin, on obtient un système triangulaire supérieur $U\vec{x} = \vec{b}^{(3)}$, que l'on résout facilement (substitution rétrograde).

\subsubsection{Avantages de la décomposition $L \cdot U = A$}

\textbf{Inconvénient de la méthode de Gauss simple :} Si l'on doit résoudre $A\vec{x}^k = \vec{b}^k$ avec de multiples vecteurs $\vec{b}^k$ qui dépendent de $k$, on doit tout recommencer à chaque fois.

\textbf{La solution avec LU :} Si on définit formellement $L \cdot U = A$, on a maintenant 2 systèmes triangulaires à résoudre à chaque étape $k$ :
\[
\begin{cases}
L\vec{y}^k = \vec{b}^k & \text{(Système triangulaire inférieur)} \\
U\vec{x}^k = \vec{y}^k & \text{(Système triangulaire supérieur)}
\end{cases}
\]

\textbf{Avantages :}
\begin{enumerate}
    \item Calculer \textbf{une fois pour toutes} la décomposition $LU$ de $A$.
    \item À chaque étape, on résout très rapidement les deux systèmes $L\vec{y}^k = \vec{b}^k$ et $U\vec{x}^k = \vec{y}^k$.
\end{enumerate}

\subsubsection{Pivotage partiel}

\textbf{Exemple problématique :}
\[ 
A = \begin{pmatrix} 1 & 2 & 3 \\ 2 & 4 & 5 \\ 7 & 8 & 9 \end{pmatrix} \xrightarrow{L_2 \leftarrow L_2 - 2L_1} \tilde{A}^{(2)} = \begin{pmatrix} 1 & 2 & 3 \\ 0 & 0 & -1 \\ \dots & \dots & \dots \end{pmatrix}
\]
Ici, le pivot devient nul ($a_{22} = 0$). L'algorithme bloque. On peut (et on doit) échanger les lignes. \\
Si on fait la même chose pour la matrice $L$, on obtient la résolution via une matrice de permutation $P$ :
\[ \tilde{L}\tilde{U} = PA \]

\textbf{Résoudre $A\vec{x} = \vec{b}$} devient équivalent à :
\[ PA\vec{x} = P\vec{b} \iff \tilde{L}\tilde{U}\vec{x} = P\vec{b} \]
On résout alors : $\tilde{L}\vec{y} = P\vec{b}$ puis $\tilde{U}\vec{x} = \vec{y}$.

\textit{En pratique, pour des raisons de stabilité numérique, on va toujours échanger les lignes de sorte à avoir le plus grand pivot possible en valeur absolue.}

\begin{algorithme}[LU avec changement de pivot]
\textbf{Entrées :} Matrice $A$, vecteur $\vec{b}$ \\
\textbf{Initialisation :} $A^{(1)} = A$, $P = I_{n \times n}$, $L = 0_{n \times n}$ \\
\textbf{Pour} $k = 1, \dots, n-1$ \textbf{faire :}
\begin{itemize}
    \item Trouver $i \ge k$ tel que $|a_{ik}^{(k)}|$ soit maximal.
    \item Échanger les lignes $i$ et $k$ de la matrice $A^{(k)}$, de $L$, et de $P$.
    \item \textbf{Pour} $i = k+1, \dots, n$ \textbf{faire :}
    \begin{itemize}
        \item $l_{ik} = a_{ik}^{(k)} / a_{kk}^{(k)}$
        \item \textbf{Pour} $j = k+1, \dots, n$ \textbf{faire :}
        \begin{itemize}
            \item $a_{ij}^{(k+1)} = a_{ij}^{(k)} - l_{ik} \cdot a_{kj}^{(k)}$
        \end{itemize}
    \end{itemize}
\end{itemize}
\textbf{Fin Pour} \\
On complète la diagonale : $l_{ii} = 1$ pour tout $i$. \\
$U$ est la partie triangulaire supérieure de $A^{(n)}$. \\
\textbf{Résolution :}
\begin{enumerate}
    \item Résoudre $L\vec{y} = P\vec{b}$ (Substitution progressive)
    \item Résoudre $U\vec{x} = \vec{y}$ (Substitution rétrograde)
\end{enumerate}
\textbf{Coût computationnel :} $\Theta(n^3)$
\end{algorithme}

\subsection{Matrice Symétrique Définie Positive (SDP)}

\begin{definition}[Matrice SDP]
Une matrice $A \in \mathbb{R}^{n \times n}$ est dite Symétrique Définie Positive (SDP) si :
\begin{enumerate}
    \item Elle est symétrique : $A^T = A$
    \item Pour tout vecteur non nul $\vec{x} \in \mathbb{R}^n \setminus \{\vec{0}\}$, on a : $\vec{x}^T A \vec{x} > 0$
\end{enumerate}
\end{definition}

\begin{theoreme}[Factorisation de Cholesky]
Soit $A$ une matrice symétrique définie positive. Alors il existe une unique matrice $L$ triangulaire inférieure à valeurs diagonales strictement positives telle que :
\[ A = LL^T \]
\end{theoreme}

Si on veut résoudre $A\vec{x} = \vec{b}$ pour une matrice SDP, on calcule sa factorisation de Cholesky, puis on résout en cascade :
\[
\begin{cases}
L\vec{y} = \vec{b} \\
L^T\vec{x} = \vec{y}
\end{cases}
\]

\subsection{Un exemple d'application : Les moindres carrés}

On considère un système surdéterminé d'équations (plus d'équations que d'inconnues) :
\[
\begin{cases}
a_{11}x_1 + \dots + a_{1n}x_n &= b_1 \\
\vdots \\
a_{m1}x_1 + \dots + a_{mn}x_n &= b_m
\end{cases}
\]
Un tel système, noté matriciellement $A\vec{x} = \vec{b}$ avec $A \in \mathbb{R}^{m \times n}$ et $m > n$, n'a généralement \textbf{pas de solution exacte}.

\textbf{Exemple concret :} On cherche à déterminer la position $x(t)$ d'un mobile au cours du temps, en supposant que le modèle est affine $x(t) = vt + x_0$. Si l'on effectue 4 mesures à des instants $t_i$, le système risque fort de ne pas posséder de droite passant exactement par tous les points.

Mais on peut s'intéresser à trouver la "meilleure approximation". On cherche à minimiser l'erreur quadratique, c'est-à-dire trouver $\vec{x}$ qui minimise la norme de l'erreur $\vec{e}$ :
\[ \min_{\vec{x} \in \mathbb{R}^n} \|\vec{e}\|_2 = \min_{\vec{x} \in \mathbb{R}^n} \|A\vec{x} - \vec{b}\|_2 \]
où la norme $L_2$ est définie par $\|\vec{v}\|_2 = \left(\sum_{i} v_i^2\right)^{\frac{1}{2}}$. On dit qu'on cherche une solution au sens des \textbf{moindres carrés}.

\begin{theoreme}[Équations normales]
Soit $A \in \mathbb{R}^{m \times n}$ avec un rang maximal ($\text{rang}(A) = n$, c'est-à-dire que ses colonnes sont linéairement indépendantes ou que $A$ représente une application injective). \\
Alors il existe un unique vecteur $\vec{x} \in \mathbb{R}^n$ qui minimise la norme $\|A\vec{x} - \vec{b}\|_2$, et ce vecteur est donné par l'unique solution du système carré régulier suivant (appelé équations normales) :
\[ A^T A \vec{x} = A^T \vec{b} \]
\end{theoreme}